\subsection*{ADR-16 --- Núcleo modular, Boletería, Operación y Trabajador; Pagos y datos compartidos}
\textbf{Estado:} decisión acordada, implementada en
\texttt{feature/}\allowbreak\texttt{sebastian-arquitectura} e integrada en \texttt{develop}
mediante el PR \#44 (2/11/2026). El corte del entorno habitual sigue pendiente: integrar
no equivale a desplegar en ese entorno.

\textbf{Contexto.} El coste de coordinar programas, versiones, llamadas y datos
superaba la separación necesaria para migrar lo ya construido. Las diez aplicaciones
existentes y los catorce servicios propuestos por ADR-14 no coincidían. El objetivo
es conservar comportamiento y dejar ubicaciones claras para que cada integrante
termine sus ocho CU; no completar sus pendientes durante esta reorganización.

\textbf{Decisión.} Un workspace backend y cuatro responsabilidades de programa:
Núcleo modular, Boletería, Operación y Trabajador. Existen los primeros dos y el
consumidor de correo del Trabajador; Operación no tiene backend terminado. Pagos
es módulo reutilizable, sin servidor propio; webhook único en Núcleo. Base
PostgreSQL única por ocho esquemas y libro financiero único, preparación de datos
coordinada antes de aplicaciones, conexiones nombradas y synchronize=false.
Se conservan contratos públicos, autorización, migraciones y autoridades de Entradas.

\textbf{Sustituye.} ADR-14 en granularidad/extracción y topología de catorce
servicios; ADR-12 en la composición de su agregado; ADR-01 en servidor/base por
dominio; ADR-06 en bases independientes/MongoDB como destino de esta migración.
Revisa parcialmente ADR-15: llamadas internas a Pagos/Ventas pasan a interfaces
locales y su webhook llega a Núcleo. Conserva gateway, RS256, roles, turno único,
HMAC y los outbox/proyecciones existentes; no convierte sus garantías propuestas
en funciones terminadas. ADR-02/03/05/07--11 conservan sus mecanismos aplicables.
Los textos y diagramas anteriores quedan identificados como historia.

\textbf{Concesiones y revisión de requisitos arquitectónicos.}
\begin{itemize}
\item Se acepta despliegue/versionado/migraciones coordinados y un dominio de fallo
SQL compartido; no independencia de base por microservicio. Los pools siguen
separados y requieren dimensionado conjunto. No se afirma aislamiento de capacidad.
\item RNF-10 conserva escalado independiente de Boletería frente al Núcleo; el
escalado de cada componente de Entradas deja de ser individual. Requiere nueva
medición, no reutilizar cifras de ADR-12/14 como cumplimiento.
\item El aislamiento de proceso de Operación se conserva como objetivo para
CU-010/019/020. La activación con Núcleo caído y alertas con acuse siguen pendientes.
\item No se incorpora MongoDB, índice de búsqueda, Kubernetes o réplica de lectura
para cerrar la migración. Sus propuestas/ASR anteriores requieren revisión del
equipo; reportes y búsquedas ausentes no se declaran cumplidos con PostgreSQL vacío.
\item Compartir SQL no aporta atomicidad entre operaciones existentes ni elimina
ventanas QR/proyecciones. F12/F13/F14/F15 continúan abiertos. No se modifica el SRS,
los umbrales RNF o el reparto para aparentar cumplimiento.
\end{itemize}

\textbf{Alternativas.} Conservar todos los servidores mantiene complejidad de red y
versiones sin evidencia de necesidad. Un único backend reduciría más procesos,
pero perdería el aislamiento previsto de Boletería/Operación. Una plataforma nueva
de trabajos o reescribir todos los CU ampliaría el alcance sin preservar primero
la referencia. Se elige migración incremental y pruebas reales sobre copias.

\textbf{Verificación y límites.} La evidencia local de \texttt{ESTADO-MIGRACION}
apartado 23 cubre 830 unitarias y recorridos transferido/vacío, 73 Núcleo, 23
Boletería y contratos financieros/correo, imágenes y gateway. Las ejecuciones
remotas, digests y bloqueos se registran en el apartado 24 y en la solicitud de
integración, sin equiparar YAML válido a CI/CD ejecutado. Corte/reversión habitual
están documentados, no ejecutados. Después de escrituras nuevas volver de imagen
no revierte datos.

\textbf{Consecuencia para el equipo.} Los responsables del SRS conservan sus CU.
Guía por integrante: \texttt{GUIA-EQUIPO-MIGRACION.md}. Completar los CU ausentes,
F12 o evacuaciones no es requisito para terminar la migración del código existente.
